\chapter{Conformance with React}\label{chap:rules}
\Rtrace{} is a faithful model of Hooks backed by both theoretical and empirical evaluation:
\begin{enumerate}
  \item \textbf{Conformance according to the official documentation:}
    Although React does not come with a formal definition, we can extract key properties of React based on the documentation.
    We show that \Rtrace{} respects the key properties of React as documented in~\cref{sec:key}.
  \item \textbf{Direct testing against the implementation:}
    We present the results of empirical tests comparing \Rtrace{} against (multiple releases of) React's runtime behavior across a range of scenarios, covering all evaluation rules, in~\cref{sec:tests}.
\end{enumerate}


\section{Conformance Theorems on React}\label{sec:key}
The React documentation provides an informal explanation of the behavioral properties of React applications.
Most of these properties are explicitly captured in our semantics.
To demonstrate that our semantics align with the behaviors described in the documentation, we prove two key properties in~\cref{subsec:key}:
\begin{description}
  \item[\Cref{thm:reeval}:] When a setter function is called during rendering, React immediately attempts to re-evaluate the component body with the updated state \citep{Met25UseState}.
  \item[\Cref{thm:eff}:] A view's Effects are executed after re-rendering triggered by a state update from itself or one of its ancestors \citep{Met25Render}.
\end{description}

While there are slight differences between our semantics and the React runtime due to an optimization in React, this does not affect reasoning about most React applications.
We have formalized this optimization and proved that it preserves output behavior (\cref{thm:simtrans},~\cref{subsec:opt}), provided that applications adhere to the constraints specified by React.


\subsection{Key Properties of React}\label{subsec:key}
The first key property states that when a setter function is called during rendering, React immediately attempts to re-evaluate the component body with the updated state.
The pitfall discussed in~\cref{subsec:topset} is due to this property.
This behavior is modeled by allowing multiple evaluations of the component body to occur within a retrying evaluation, which happens iff a setter is applied during the first evaluation.

\begin{thm}[Re-Evaluation by Calling Setter]\label{thm:reeval}
The derivation of \(\Step*[\pi, \sigma]{e}{v', \pi'', \omega'}\) includes multiple evaluations of \(e\)
iff \Rule{AppSetComp} appears in the derivation of the first evaluation:
\[ \Step[\pi\{\fv{dec}{\pi.\field{dec} \setminus \{\dec{Check}\}}, \fv{effq}{[]}\}, \sigma]{e}{v, \pi', \omega} \]
\end{thm}

\begin{proof}
Multiple evaluations of \(e\) occur in a retrying evaluation iff the resulting view~\(\pi'\) from the first evaluation includes the \dec{Check} decision (\Rule{EvalMult}).
Thus, it suffices to show that \(\pi'.\field{dec}\) contains \dec{Check} iff \Rule{AppSetComp} appears in the derivation of the first evaluation.

This follows directly from the semantics: \Rule{AppSetComp} is the only rule that introduces \dec{Check} during evaluation (when \(\phi \in \{\phase{Init}, \phase{Succ}\}\)).
Moreover, once added, no rule removes \dec{Check} from the view's decision.
Therefore, \(\pi'\) includes \dec{Check} iff a setter is applied during the first evaluation.
\end{proof}

The second key property states that Effects are executed when queued state updates modify the state values of a view or one of its ancestors (case~1 of \cref{thm:eff}).
Notably, Effects may still be executed even if the final state values remain unchanged, provided state setters are invoked during the evaluation of the component body (case~2 of \cref{thm:eff}).

\begin{thm}[Effect Execution Condition]\label{thm:eff}
If $\<t,m,\omega,\delta,\checkmode\> \smallstep \<t,m',\omega',\delta,\mu\>$,
then for all paths~$p$ reachable from the root in both~$m$ and~$m'$,
\Rule{CommitEffsPath} with~$p$ appears in the derivation of the next transition~$\<t,m',\omega',\delta,\mu\> \smallstep \<t,m'',\omega'',\delta,\mu'\>$
iff there exists an ancestor~$p'$ of~$p$ such that either
\begin{enumerate}
  \item for some $\ell \in \domain(m'[p'].\field{sttst})$, we have $m[p'].\field{sttst}[\ell].\field{val} \ne m'[p'].\field{sttst}[\ell].\field{val}$; or
  \item the derivation of $\<t,m,\omega,\delta,\checkmode\> \smallstep \<t,m',\omega',\delta,\mu\>$ includes \Rule{AppSetComp} with path~$p'$.
\end{enumerate}
\end{thm}

\begin{proof}
We first show only the views in rendered mode~\(\rendermode\) may carry the \dec{Effect} decision. This follows from the operational rules governing state transitions:
\begin{itemize}
  \item \Rule{StepInit} transitions the initial state to rendered mode~\(\rendermode\).
  \item \Rule{StepEffect} clears all \dec{Effect} decisions from descendants of~$t$ via \(\sem{commitEffs}\).
  \item \Rule{StepCheck} ensures \dec{Effect} is added only when transitioning into rendered mode~\(\rendermode\), by induction on \(\sem{check}\).
  \item \Rule{StepEvent} does not introduce \dec{Effect}.
\end{itemize}

The \Rule{CommitEffsPath} rule with path~$p$ appears in the derivation of the next transition iff the \Rule{StepEffect} rule is applied (i.e., \(\mu = \rendermode\)) and \dec{Effect} is included in $m'[p].\field{dec}$.
Among the views reachable in~$m'$, only those updated by \(\sem{check}\) during the first transition may carry the \dec{Effect} decision.
More precisely, \dec{Effect} is added to~$p$ iff~$p$ is a descendant of some~$p'$ such that \(\sem{check}(p')\) is derived via \Rule{CheckEffect}.
This holds iff the resulting view~\(\pi'\) from the retrying evaluation of~$p'$'s body satisfies \(\dec{Effect} \in \pi'.\field{dec}\).

We now show that \dec{Effect} is added to such a view after the retrying evaluation iff either (a)~the final state differs from the initial state, or (b)~setters are applied during the evaluation.

Case~(a) is immediate: state updates occur only via \Rule{SttReBind}, which adds \dec{Effect} when the new value differs from the old.

In case~(b), applying a setter adds \dec{Check} and triggers re-evaluation via \Rule{EvalMult}.
Resolving this re-evaluation loop requires the subsequent evaluation to yield a different result, which necessitates a different binding via \Rule{SttReBind}, thereby introducing \dec{Effect}.
Hence, any terminating derivation involving \Rule{AppSetComp} includes \dec{Effect} in the resulting decision.

Conversely, suppose \dec{Effect} is added, but the final state values are equal to the initial ones.
Then at least one state must have changed and reverted---that is, some label~$\ell$ was updated from~$v$ to~$v'$, and then back to~$v$---implying multiple applications of \Rule{SttReBind} with the same label.
This only occurs when \Rule{AppSetComp} triggers multiple evaluations of the body.
Therefore, \dec{Effect} is added iff either some state value differs from its original value, or \Rule{AppSetComp} appears in the derivation.
\end{proof}


\subsection{Optimization in React}\label{subsec:opt}
There is a subtle discrepancy between our semantics and the React runtime due to an optimization:
When a state update callback returns the same value as the current state, React skips re-evaluating the component body.

\begin{wrapfigure}[5]{l}{0.43\linewidth}
  \vspace*{-1.4\intextsep}
  \begin{reactcode}
let Counter _ =
  print 0;
  let (s, setS) = useState 1 in
  [s, button (fun _ ->
   setS (fun s -> (print 1; s));
   setS (fun s -> (print 2; s+1));
   setS (fun s -> (print 3; s+1)))];;
Counter ()
  \end{reactcode}
\end{wrapfigure}
\noindent According to our semantics, clicking the button prints~\verb|0\0\1\2\3|.
The first~\verb+0+ is printed during the initial render.
When the button is clicked after the render, the event handler queues the state updates without printing anything.
Then, during the second render triggered by the state updates, another~\verb+0+ is printed.
Finally, as the updates are applied,~\verb+1\2\3+ are printed.

In contrast, React prints~\verb+0\1\2\0\3+ due to the optimization.
The first~\verb+0+ is printed during the initial render.
Upon clicking the button, the first two updates are evaluated immediately:~\verb+1+ and~\verb+2+ are printed as React checks whether the state has changed.
The second update returns a new state, triggering a re-render.
The third update is skipped at this point because a re-render is already scheduled.
During the re-render, the second~\verb+0+ is printed, followed by~\verb+3+ from the queued third update.

We could have included this optimization for completeness, but we chose not to in order to maintain conciseness.
However, this omission does not compromise correctness as long as applications conform to the constraints specified in the React documentation.

% To model the optimization formally, we introduce the notions of purity, normalization, and similarity,
% and
Indeed, we show that the optimization preserves program behavior in \cref{thm:simtrans} (under certain conditions).

The optimization can be understood as partially applying queued state updates without re-executing the component body.
Assuming these updates are pure (\cref{def:purity}) as required by the React documentation, this corresponds to a form of partial normalization,
in which only a prefix of the pure updates is applied.
We formalize the relationship between the original and optimized tree memory using the notion of \emph{similarity} ($m \approx_t m'$; \cref{def:similarity}),
which is defined in terms of \emph{normalization} (\cref{def:normalization})---a process of applying pure updates in advance.

\begin{defn}[Purity]\label{def:purity}
  A closure~$\<\lfun x e, \sigma\>$ is \emph{pure} iff
  the application of any value~$v$ does not cause any side effects.
  That is, for all~$\Sigma$ and~$v$, we have~$\Step[\Sigma,\sigma[x\mapsto v]]{e}{\_,\Sigma, []}$.
\end{defn}

\begin{defn}[Normalization]\label{def:normalization}
  A \emph{normalization} of a state store entry is defined as the result of applying its pure prefix of state updates, with its resulting decision collected.
  Let~$l$ be the length of such a prefix.
  Then
  \begin{multline*}
    \sem{normalize}\bigl(\bigl\{\fv{val}{v_0},\fv{sttq}{\bigl[\Overline{\<\lfun{x_i}{e_i},\sigma_i\>}\bigr]_{i=1}^n}\bigr\}\bigr)
    = \bigl\<\bigl\{ \fv{val}{v_l},
         \fv{sttq}{\bigl[\Overline{\<\lfun{x_i}{e_i},\sigma_i\>}\bigr]_{i=l+1}^n} \bigr\}, d\bigr\> \\
    \text{where }\bigl(
      \Step[\pi,\sigma_i[x_i \mapsto v_{i-1}]]{e_i}{v_i,\pi, []}
    \bigr)_{i=1}^{l}\\
    \text{and }d =
      \scondp{l \ne n \scondq \{\dec{Check}\}
      \scondcolon v_0 \nequiv v_l \scondq \{\dec{Check},\dec{Effect}\}
      \scondcolon \{\}}
  \end{multline*}

  We extend this notion to a state store~$\rho$ and a view~$\pi$, i.e., $\sem{normalize}(\rho)$ normalizing all of its entries, and~$\sem{normalize}(\pi)$ normalizing its state store:
  \begin{flalign*}
    \hspace{69.9pt}\sem{normalize}(\rho) &= \begin{multlined}[t]\Bigl\<
        [\Overline{\ell \mapsto r_\ell}]_{\ell \in \domain \rho},
        \smashoperator{\bigcup_{\ell \in \domain \rho}} d_\ell
    \Bigr\> \\ \text{where } \sem{normalize}(\rho[\ell]) = \<r_\ell,d_\ell\>,
    \end{multlined}\\
    \sem{normalize}(\pi) &= \begin{multlined}[t]\pi \bigl\{ \fv{sttst}{\rho'},\fv{dec}{\pi.\field{dec} \setminus \{\dec{Check}\} \cup d} \bigr\} \\ \text{where } \<\rho',d\> = \sem{normalize}(\pi.\field{sttst}).\hspace{70pt}\qedhere % TODO: qedhere does not work here well
    \end{multlined}
  \end{flalign*}
\end{defn}

\begin{defn}[Similarity]\label{def:similarity}
  Two views are \emph{similar} iff they are equal under normalization:
  \[
    \pi \approx \pi'
    \triangleiff
    \sem{normalize}(\pi) = \sem{normalize}(\pi')
  \]

  We extend this notion to memories, i.e., $m$~and~$m'$ are \emph{$t$-similar} iff the descendants of~$t$ in $m$~and~$m'$ are all similar and the rest of them are equal:
  \[
    m \approx_t m'
    \triangleiff
    {}\land \begin{lgathered}
      \forall p \in \sem{reachable}(m, t),\ m[p] \approx m'[p] \\
      \forall p \notin \sem{reachable}(m, t),\ m[p] = m'[p] \\
    \end{lgathered}
  \]
  where $\sem{reachable}(m, t)$ is the set of all paths in~$m$ that is reachable from the tree~$t$.
\end{defn}

We now state that evaluating a view similar to the original yields the same result (\cref{lem:simevalbody}), provided the view is \emph{valid} (\cref{def:validity}).

\begin{defn}[Validity]\label{def:validity}
  A view is \emph{valid} if it has only the states that are present in its component body.
  More precisely,~$\pi$ is \emph{valid under~$\delta$} iff the domain of $\pi$'s state store is exactly the set of state labels of~$e$, where~$\pi.\field{spec} = \<C, v\>$ and~$\delta[C] = \<\lfun{x}{e}, \sigma\>$.
  We extend this notion to tree memory: A tree memory~$m$ is valid under~$\delta$ iff its views are all valid under~$\delta$.
\end{defn}

Note that all views encountered during execution are valid.
This follows from a syntactic restriction on component bodies: Hooks must appear only at the top level.
This ensures that each \reacttt{useState} call in a component body is executed exactly once during initialization,
so all views produced by~$\sem{init}$ are valid.

\begin{restatable}[Similar Evaluations of Component Body]{lem}{simevallem}\label{lem:simevalbody}
  Evaluating the component body of similar views produces the same value, view, and output buffer.
  That is, for $\pi$ valid under $\delta$ where $\pi.\field{spec} = \<C,v\>$ and $\delta[C] = \<\lambda x.e,\sigma\>$,
  if $\Step*<\phase{Succ}>[\pi,\sigma[x\mapsto v]]{e}{v',\pi', \omega}$,
  we have $\Step*<\phase{Succ}>[\hat \pi,\sigma[x\mapsto v]]{e}{v',\pi', \omega}$ for all~$\hat\pi \approx \pi$.
\end{restatable}

Building on \cref{lem:simevalbody}, we conclude that the optimization---which replaces some views in memory with similar ones---preserves program behavior.
Note that the output buffer~$\omega$ may differ if component bodies perform printing,
since the optimization may omit their evaluation entirely.
Nevertheless, the resulting memory remains identical, regardless of whether the optimization is applied.

\begin{restatable}[Similar Transitions]{thm}{simtransthm}\label{thm:simtrans}
  If a program transitions to a state in check mode~$\checkmode$ during execution, replacing it with a similar state results in the same final state as the original transition.
  That is, if $\<e,\delta\> \smallstep^* \<t,m,\omega,\delta,\checkmode\> \smallstep \<t,m',\omega',\delta,\mu\>$, then for any $\hat m$ such that $m \approx_t \hat m$, we have $\<t,\hat m,\omega,\delta,\checkmode\> \smallstep \<t,m',\omega'',\delta,\mu\>$.
  We also have~$\omega' = \omega''$ when the component bodies do not print.
\end{restatable}

Full proofs are provided in~\cref{app:chap:proofs}.

\section{Conformance Test Suite}\label{sec:tests}
Our test suite, which covers all~44 evaluation rules\footnote{See~\cref{app:sec:sem} for the complete set of rules.} comprising the operational semantics~(\cref{sec:sem}), contains 38~tests that cover 18~scenarios shown in \cref{tab:test} that compare the behaviors of programs under \Rtrace{} against React's behavior.
We have constructed the tests ourselves, testing the \Rtrace{} interpreter (\cref{chap:impl}) as well.

\begin{table}[t]
  \caption{Empirical validation of \Rtrace{} against React behavior.}\label{tab:test}
  \centering\smaller[1]
  \begin{tabular}{rlrc}
    \toprule
        & \textbf{Scenarios} & \textbf{Tests \#} & \textbf{\Rtrace{}} \\
    \midrule
     S1 & No re-render w/o a setter call & 6 & \checkmark \\
     S2 & Retries ($0<n<25$) w/ setter call during body eval & 4 & \checkmark \\
     S3 & Infinite retries ($n\ge25$) w/ setter call during body eval & 1 & \checkmark \\
     S4 & No re-render w/o Effects w/ setter call during body eval & 1 & \checkmark \\
     S5 & No re-render w/ Effect w/o setter call & 2 & \checkmark \\
     S6 & No re-render w/ Effect w/ id setter call & 1 & \checkmark \\
     S7 & No re-render w/ Effect w/ setter calls composing to id & 2 & \checkmark \\
     S8 & Re-renders ($0<n<100$) w/ Effect w/ setter call & 16 & \checkmark \\
     S9 & Infinite re-renders ($n\ge100$) w/ Effect w/ diverging setter call & 2 & \checkmark \\
    S10 & Re-render w/ child updating parent during Effect & 2 & \checkmark \\
    S11 & Re-render w/ sibling updating another during Effect & 1 & \checkmark \\
    S12 & Error w/ child updating parent during body eval & 1 & \checkmark \\
    S13 & Non-trivial reconciliation & 4 & \checkmark \\
    S14 & No re-render w/ direct object update & 1 & \checkmark \\
    S15 & Re-render w/ idle but parent updates & 2 & \checkmark \\
    S16 & User event sequence & 6 & \checkmark \\
    S17 & Re-render w/ setter call from user event & 4 & \checkmark\rlap{${}^\dagger$} \\
    S18 & Recursive view hierarchy & 2 & \checkmark \\
    \bottomrule
    && $38\mathrlap{{}^*}$ & \checkmark \\
    &&&\hfill\llap{\smaller[1]${}^*$Some tests cover multiple scenarios.}\\
    &&&\hfill\llap{\smaller[1]${}^\dagger$React's optimization changes some execution orders.}
  \end{tabular}
\end{table}

It is difficult to unit test a single evaluation rule independently because running a React program inevitably leaves footprints on multiple rules, and we have chosen the test scenarios to exhibit the pitfalls explained in~\cref{chap:pit}, as well as to check the core aspects of React including render counts, side effect ordering, reconciliation, and event handling.
For instance, although one of our test cases\footnote{\texttt{effect\_queue\_gets\_flushed\_on\_retry}} checks whether the effect queue gets flushed on retry, effectively testing rule \Rule{EvalMult}~(\cref{fig:evalmult}), it is simply categorized as S2~and~S8 in \cref{tab:test}.

Notably, we test the boundaries of React that even experienced developers may find confusing.
For instance, S12~in \cref{tab:test} covers attempting to update a parent's state during the evaluation of its child, which correctly produces an error in both implementations.
Other scenarios include unnecessary re-rendering that modifies the view hierarchy~(S8--S11,~\cref{tab:test}; similar to examples in~\cref{sec:unnec}), constructing complex UI with recursive components~(S18,~\cref{tab:test}), and more.

The tests consist of equivalent components implemented twice---once in \textsc{\textsf{React\--tRace}} and once in React---then comparing that they exhibit the same behavior.
For example, to empirically check infinite retries~(S3,~\cref{tab:test}) on the React-side, we install an error boundary \citep{Met25Component} to catch an exception when a component reaches React's hard-coded limit of 25~retries.
On the \Rtrace{}-side, we set the retry threshold to~25 and see whether the component did not stop rendering.
For tests measuring the number of render cycles on the React-side, we indirectly measure the count by counting the prints inside an Effect.

Almost all tests show identical behavior between \Rtrace{} and React's runtime, with one minor difference due to the optimization performed by React as discussed in~\cref{subsec:opt}.
We might as well include this optimization in \Rtrace{} by adding a special flag and modifying \Rule{AppSetNormal}, but we deliberately chose not to in order to keep the clarity of our semantics.

\Rtrace{} does not model a specific version of React: We have tested against the latest versions of every major React release since Hooks were introduced in~\href{https://www.npmjs.com/package/react/v/16.8.0}{16.8 (Feb~2019)}.  % https://legacy.reactjs.org/docs/hooks-reference.html
All test cases have been reproduced in versions~\href{https://www.npmjs.com/package/react/v/16.14.0}{16.14.0\,(Oct~2020)}, \href{https://www.npmjs.com/package/react/v/17.0.2}{17.0.2 (Mar~2021)}, \href{https://www.npmjs.com/package/react/v/18.3.1}{18.3.1\,(Apr~2024)}, and \href{https://www.npmjs.com/package/react/v/19.1.0}{19.1.0\,(Mar~2025)}.

The test suite is available at \href{https://github.com/Zeta611/react-trace}{\gh{Zeta611/react-trace}} alongside the \Rtrace{} interpreter.
